\textbf{\textit{Soal.}} Let $ABC$ be a triangle. Let $L$ be the line through $B$ perpendicular to $AB$. The perpendicular from $A$ to $BC$ meets $L$ at the point $D$. The perpendicular bisector of $BC$ meets $L$ at the point $P$. Let $E$ be the foot of the perpendicular from $D$ to $AC$.

Prove that triangle $BPE$ is isosceles.

\begin{proof}[\textbf{Solusi. }] (Thursday, 24 July 2025) 
    Misalkan garis sumbu $BC$ memotong $BC$ dan $BE$ berturut-turut di $M$ dan $N$. Lalu misalkan pula $T$ sebagai perpotongan $AD$ dengan $BC$.
    Dari definisi garis sumbu mudah dilihat bahwa $BN=NC$, $BP=PC$, $\dangle NPB = \dangle CPN$,  $\dangle BCP = \dangle PBC$. 
    Perhatikan bahwa $\dangle ABD = \dangle AED = 90^\circ$ yang menyebabkan $ABDE$ siklis.
    Lalu, kita juga punya $AD \perp BC$ dan $NP \perp BC$ sehingga $AD \parallel NP$.
    Oleh karena itu, 
    \begin{align*}
        \dangle CEN = \dangle AEB = \dangle ADB = \dangle NPB = \dangle CPN
    \end{align*}
    yang menyebabkan $NECP$ siklis.
    \begin{center}
    \begin{asy}
        import olympiad;
        import geometry;
        // Set the size of the diagram
        size(250);

        // Define the vertices of triangle ABC
        pair A, B, C;
        B = (0,0);
        A = (1,6);
        C = (8,1);

        // Define the direction of Line L (perpendicular to AB)
        pair L_dir = rotate(90)*(A-B);
        // Draw the line L extending in both directions from B
        path L = (B - 7*unit(L_dir)) -- (B + 10*unit(L_dir));

        // The line from A perpendicular to BC
        pair A_perp_foot = foot(A, B, C);
        path line_AD = A--A_perp_foot;

        // Point D is the intersection of L and the perpendicular from A to BC
        pair D = extension(B, B + L_dir, A, A_perp_foot);

        // The perpendicular bisector of BC
        pair M = (B+C)/2; // Midpoint of BC
        pair perp_bisector_dir = rotate(90)*(C-B);
        path perp_bisector = M -- M + 5*unit(perp_bisector_dir);

        // Point P is the intersection of L and the perpendicular bisector of BC
        pair P = extension(B, D, M, M + perp_bisector_dir);

        // Point E is the foot of the perpendicular from D to AC
        pair E = foot(D, A, C);

        // Point H is the intersection of BE and the perpendicular bisector of BC
        pair N = extension(B, E, M, M + perp_bisector_dir);

        // T intersection of AD and BC
        pair T = intersectionpoint(line(B,C),line(A,D));

        // --- Drawing Commands ---

        // Draw the main triangle ABC
        draw(A--B--C--A);

        // Draw the construction lines
        draw(A--D, red);
        draw(M--P--N, darkgreen);

        // Draw the triangle BPE, which is to be proved isosceles
        draw(B--P--E--P--C, magenta);
        draw(B--E);
        draw(N--C);

        // Add dots and labels for all key points
        dot("$A$", A, NW);
        dot("$B$", B, SW);
        dot("$C$", C, SE);
        dot("$D$", D, SW);
        dot("$P$", P, SW);
        dot("$E$", E, NE);
        dot("$M$", M, SW);
        dot("$N$", N, NW);
        dot("$T$", T, NE);

        // Add right angle markers to show perpendicular lines
        rightanglemark(A, B, D, 5);
        rightanglemark(A, A_perp_foot, C, 5);
        rightanglemark(D, E, A, 5);
        rightanglemark(P, M, B, 5);
    \end{asy}
\end{center}
    Di lain sisi, karena $\dangle TAB = \dangle DAB$ dan $\dangle DBA = 90^\circ = \dangle ATB$, maka $\dangle CBA = \dangle ADB$, yang menyebabkan $\triangle AEB \sim \triangle ABC$.
    Dari sini didapat $\dangle ABE = \dangle ECB$. 
    Dari fakta-fakta tersebut kita punya
    \begin{align*}
        \dangle PEC
        &= \dangle PNC \\
        &= \dangle BNP \\
        &= 90^\circ - \dangle CBN \\
        &= \dangle DBA - \dangle CBN \\
        &= \dangle EBA + \dangle PBC \\
        &= \dangle ECB + \dangle BCP \\
        \dangle PED &= \dangle ECP.
    \end{align*}
    Berarti diperoleh $EP=PC$. Namun, karena $BP=PC$ maka $EP=BP$. Terbukti bahwa $\triangle BEP$ sama kaki.
\end{proof}